% introduccion

A lo largo de este informe analizamos el flujo de los paquetes en tres redes ubicadas en distintos lugares del área metropolitana de Buenos Aires. Pudimos observar que la mayoría de los paquetes son de tipo \textit{$\langle$Unicast, IPv4$\rangle$} y que varía fuertemente el caudal de paquetes cuando la red está en uso por el host y cuando no. 

A su vez, observamos que el porcentaje de paquetes de tipo \textit{Broadcast} es prácticamente despreciable en comparación con los de tipo \textit{Unicast}. Esto tiene sentido, ya que los paquetes del mismo tipo suelen ocurrir en el contexto de protocolos de control.

Sin embargo, consideramos que las muestras tomadas no fueron lo suficientemente representativas del comportamiento de la red. El momento del día, el tipo de conexión, la ubicación geográfica y la actividad del host son distintos factores que definen su comportamiento. La cantidad de muestras que tomamos no son suficientes para capturar estos aspectos correctamente: capturamos momentos aislados, sobre acceso específicos y en contextos particulares.

Finalmente, cabe destacar que la entropía máxima para una fuente de información nula se da en el caso en que los simbolos ocurren con equiprobabilidad, y  corresponde con $\log_2(n)$ ---donde $n$ representa la cantidad de símbolos presentes---. Notamos que ninguna de las fuentes tuvo una entropía que se acerque a este límite teórico. 
